\chapter{Bewegung}
Bevor die Umsetzung der Erfassung geklärt werden kann, muss definiert werden, was unter Bewegung zu verstehen ist, um die optimale Stelle am Körper für die Datenerfassung zu finden.

\section{Definition der Bewegung in Bezug auf das System}
In dieser Arbeit werden drei verschiedene Bewegungsarten analysiert: Stehen, Gehen und Rennen. Jede dieser Bewegungsarten zeichnet sich durch spezifische biomechanische Merkmale aus. Beim Gehen bewegt sich der Körper analog zu einem umgekehrten Pendel – die kinetische Energie der ersten Standphase wird in potenzielle Gravitationsenergie umgewandelt und anschließend teilweise zurückgewonnen, wenn der Körper nach vorne und unten fällt \cite{cappellini2006motor}. Im Gegensatz dazu ähnelt das Rennen dem Hüpfen auf einem Pogo-Stick, bei dem kinetische und potenzielle Gravitationsenergie vorübergehend als elastische Dehnungsenergie in Muskeln, Sehnen und Bändern gespeichert und in der zweiten Hälfte der Standphase nahezu vollständig zurückgewonnen wird \cite{cappellini2006motor}. Diese Unterschiede ist auf die Abbildung \ref{fig:walking_running} zu erkennen.

Ein weiterer bedeutender Unterschied zwischen diesen Bewegungen liegt in den Phasen des Bodenkontakts: Während das Gehen durch eine doppelte Stützphase gekennzeichnet ist, in der beide Füße kurzzeitig den Boden berühren, weist das Rennen eine „Flugphase“ auf, in der sich keine der Gliedmaßen am Boden befindet \cite{cappellini2006motor}. Diese Unterschiede sind für die spätere Klassifizierung der Bewegungen essenziell, da sie als zentrale Merkmale zur Unterscheidung genutzt werden können.

\begin{figure}[h]
	\centering
	\includegraphics[width=0.5\linewidth]{walkingrunning}
	\caption{Walking and Running  \cite{cappellini2006motor}}
	\label{fig:walking_running}
\end{figure}

\section{Welche Körperteile werden genutzt?}
Um die Unterschiede zwischen Stehen, Gehen und Rennen zu analysieren, muss man sich überlegen, welche Körperteile benutzt werden, um zu wissen, wo der Sensor platziert werden muss. Zu Beginn der Beschleunigung sind die Quadrizeps, also die Oberschenkelmuskeln, am stärksten beteiligt. Der Musculus soleus, ein Skelettmuskel, der zur Unterschenkelmuskulatur gehört, und der Musculus gastrocnemius, ein Skelettmuskel, der zur Unterschenkelmuskulatur gehört, sind die wichtigsten Muskeln. Obwohl die Arme nur wenig zur direkten Vorwärtsbewegung oder zum Abstützen beitragen, sind sie für den Ausgleich der vertikalen Drehbewegung der unteren Extremitäten von entscheidender Bedeutung. \cite{hamner2010muscle} Daher sollte der Sensor am besten an den Armen und Beinen angebracht werden. Diese möglichen Positionen sind in der Abbildung \ref{fig:sensor} dargestellt.

\begin{figure}[h]
	\centering
	\includegraphics[width=0.5\linewidth]{sensor}
	\caption{Mögliche Sensorplatzierungen (eigene Abbildung)}
	\label{fig:sensor}
\end{figure}
